% !TeX root = ../papers/nc-p3.tex
\section{平行な二重色で片方の文字対に限られる薄い関係}
\label{sec:klein-parallel-flags}
\begin{theorem}\label{thm:klein-parallel-flags}
両切断で厚い色を $0,1$ とし、高位ビットで文字を二対 $\{0,1\},\{2,3\}$ に分ける。
ある位置で、薄い色の関係の行の台が一方の文字対だけに含まれるとする。
平行方向の表を零始まりで番号付けすると、正則な台はちょうど
\[
 50,51,52,54,55,56,57,58,59,62,63,64,65,66
\]
である。これらの台では指定の二切断階数を持つ全係数が二次PBWで左右Koszul AS正則である。
さらにこの行条件を課さない部分でも、表の $76,78,112,115,117$ は実現でき、同じ全点での結論を持つ。
\end{theorem}
\begin{proof}[係数に依らない辞書順]
前の表の148代表のうち、上の行条件は0番から66番にちょうど対応する。
行集合および列集合は四位置を巡るので条件は巡回・反転で保たれる。
位置 $p$ ごとに文字順 $\pi_p$ を選び、語に位置別の左辞書順を入れる。
一色の行 $i$ に対応する二文字語は $i,i+b$ の一個なので、先頭語は行を大きい順に調べて求まる。
階数1なら最初の非零行、階数2なら最初の非零行と、それと独立な後続の最初の行を取る。
零成分が強制する同一直線の類に属する行は飛ばせる。非零成分の台が異なれば二行は必ず独立である。
同じ台だが同一直線と決まらない場合には、その順序では証明を行わない。

各行 $i$ でこの規則が選ぶ相手の集合を $F_i$ とする。$\pi_p$ の順に相手の集合が
大きさ $0,1,2,3$ の入れ子になるなら、足りない最後の文字を加えて次位置の順序 $\pi_{p+1}$ が定まる。
これを四回繰り返して元の順序に戻るものを選ぶ。四位置の $4!$ 通りの初期順序の検査で、
定理に挙げた19台に加え、$77,113,114,116$ の台がこの検査を通る。
各位置で先頭語は番号の降順対の六個であり、長さ3で残る語は20個である。
四つの一文字縮約は異なる文字次数を持ち非零なので、三次関係の交叉が次元4以上で、商の次元は20以上である。
従って全ての三文字重複が解消する。Diamond補題による非減少語の基底と二次双対性から
全次数のKoszul性を得る。$W$ の縮約による $1$ 対 $3$、$2$ 対 $2$ の対合が完全なので、
左右AS性も従う（\path{sections/unmarked_klein_pbw.tex}の辞書順の証明を参照）。
\end{proof}

\begin{proof}[実現と残る53台の全排除]
二切断の必要小行列式で、非零語の単項式を除いて二項式となるものをSmith標準形で解く。
右辺の符号を含め、指数の $1,2,4$ の全ての根を列挙し、連結する対角基底変更を割る。
各代表の係数はこうして得た自由変数のLaurent単項式と定数位相の積で与えられる。
他の小行列式を代入し、変数積を可逆化する。23台でこの解集合が空なのは
$77,113,114,116$ だけである。非空なら上記の行主元が下からも階数を保証するため、
どの解も標準AS窓である。この表示は位相を固定した係数商も再現し、例外値に追加の開条件を要求しない。

残りの0--66番の53台を排除するには、Laurent単項式の比だけで決まる二列基底を使う。
各厚い色で、ある二次小行列式がLaurent単元となる二列を固定する。薄い色では非零列を一つ固定する。
指定階数を持つ任意の実現でこれらは色関係の基底である。残る小行列式をまだ課さない、より大きい
Smithパラメータトーラス上で、固定した列に二文字を付けた長さ4・零文字の関係行列を作る。
その関数体で商次元が9を超えるなら、実現する全特殊化で同じ障害が残る。
5,13,27,36,47番の一部の位相ではこの上界だけを使わず、残る小行列式が可逆化後に単位イデアルとなることを確認する。
これで53台の全位相が尽くされる。
\path{computations/verify_klein_parallel_flags.py} は台と順序の全比較、両切断の係数正規化、
基底の単元小行列式と各位相の障害となる始点、19台の非空性を検算する。
\end{proof}
